% =====================================================================
\section{What this briefing is, and how to read it}
\label{sec:intro}
% =====================================================================

\subsection{Purpose}

This document is a technical briefing, not a journal article. It has one
reader --- the author, preparing to present, defend and extend this work
during an internship whose task is a rebalancing strategy built on the
indicator studied here --- and it is organised so that any result can be
explained from it directly: what was computed, from which formula, taken from
which reference, validated how, and trusted how far. Everything numerical in
the chapters that follow is regenerated from the pipeline's stored outputs
whenever the analysis is re-run, and every interpretive sentence is generated
from the estimates it describes, so the prose cannot drift from the tables.

\subsection{The sentence this study is built on}

\textcite{Ohta2026NoiseTraders} asks what kind of participant behaviour produces
\emph{price clustering} --- the disproportionate share of trading that executes at
round prices --- and answers it with an order-flow mechanism. Noise traders leave
limit orders at round prices outside the best quote, because working out a
precise valuation is expensive and a round number is a cheap substitute. They are
slow to revise or cancel those orders, because managing orders at speed is also
expensive. When value moves, faster participants take the stale orders with large
market orders. Three things follow, and the paper documents all three: clustering
concentrates in large trades rather than small ones, it is stronger where
individual investors are more present, and trades at round prices carry about a
basis point more price impact than trades elsewhere, because they execute against
prices that have gone stale.

The paper then closes with a claim and a gap:

\begin{quote}
\itshape
In light of these results, it may be possible to construct a measure of noise
trader activity from price clustering measures. \dots\ In that case it may be
possible to analyze the relationship between noise trader activity and price
formation and liquidity at a daily or shorter frequency. Such analysis, however,
remains a task for future work.
\end{quote}

This study is that analysis, on a stretch of tape three years past the end of
the paper's sample.

\subsection{Why the answer is not obvious in advance}

In \textcite{Kyle1985} noise traders are what makes a market liquid: they are the
volume informed traders hide inside, and without them the spread widens to
protect against adverse selection. \textcite{PeressSchmidt2020} find the
empirical counterpart --- when noise traders are distracted, liquidity gets
worse. On that reading, high clustering should mark a \emph{more} liquid market.

The mechanism Ohta documents points the other way. If clustering marks orders
that are stale and about to be picked off, it marks a market where liquidity is
being supplied by participants who are about to regret it, and where the price
impact of a large order is unusually high. On that reading, high clustering
should mark a market that is \emph{more} expensive to trade.

Both stories are about the same participants. They differ on whether what matters
is that the uninformed are present or that they are slow. Either answer would be
a result, and the sections below report which way this sample falls --- with the
robustness work deciding how hard the answer may be leaned on.

\subsection{What this study can and cannot do}

Ohta identifies noise traders using two proxies for individual-investor activity:
margin-trading balances and ownership shares. Neither is available here. That
sounds like a handicap and is partly the point: the paper's own conclusion is
that the clustering measures \emph{are} the proxy, so a study that uses only
clustering and the order book is testing the claim rather than assuming it.

Working without those proxies also forces a substitution that turns out to be an
advantage. The data here carry the ten best price levels on each side of the book
after every event. That makes it possible to see the standing round-price
inventory directly --- what share of visible resting depth is sitting at round
prices at any moment --- where a study confined to executed trades sees
round-price orders only after someone has taken them. The same data support
inferring limit-order submission and cancellation volume from how the displayed
ladder changes, which is how Ohta builds his own order-flow variables.

The limits are equally plain. Four months is not a long panel. There is no
experiment and no instrument, so nothing below is causal; the results are
descriptive associations and predictive relationships, and are labelled as such
throughout. The ten visible levels are a narrow window onto a book where most
resting volume sits further out.

\subsection{The findings at a glance}

Table~\ref{tab:findings} states every finding with the verdict the robustness
work assigns it. When presenting this study, that table is the honest summary;
each row's section gives the design and the caveats in full.

\begin{table}[htbp]
\centering
\caption{Findings and their verdicts, in one place}
\label{tab:findings}
\small
\begin{threeparttable}
\begin{tabular}{@{}p{0.30\linewidth}p{0.30\linewidth}p{0.34\linewidth}@{}}
\toprule
Finding & Headline estimate & Verdict \\
\midrule
F1 Clustering levels replicate & $M^{BLarge0}$ 14.95\% vs paper 14.8\% & Matches; validates the measurement \\
F2 Round-price impact premium & $+0.43$/$+0.51$ bp at 60s & Positive, about 34\% of the paper's one-minute size; gone by five minutes \\
F3 Size gradient & Q1 14.83\% vs Q5 15.82\% (pooled) & Absent in this screened universe \\
F4 Clustering $\to$ next-day spread & 0.0188 ($t=2.51$) & Fragile: it falls below conventional significance once the outcome's history means five lags rather than one; its sign reverses between the within-stock and cross-sectional designs \\
F5 Clustering $\to$ next-day depth & -0.0441 ($t=-5.22$) & Secondary outcome, same design and caveats \\
F6 Within-day (30-min buckets) & 0.0189 ($t=8.80$) & The robust version of the liquidity link \\
F7 Placebo on digits 1--9 & 0 of 9 significant & Result specific to the round digit \\
F8 Mechanism variables ($L$, $L^{C}$, RDepth) & RDepth $t=11.4$; $L$, $L^{C}$ null & 1 of 3 sign predictions confirmed; flow shares censored at 10 levels \\
F9 Ladder inference vs paper & $L^{S0}$ 0.102, $L^{S0C}$ 0.91 & Submission shares match (0.102 vs 0.106); cancellation ratios overshoot (0.91 vs 0.81) \\
F10 Strategy demonstration & net $-35.1$ bp/day daily; $-15.5$ smoothed & Costs dominate at daily horizon; smoothing cuts costs 32 to 14 bp/day, still negative net \\
\bottomrule
\end{tabular}
\begin{tablenotes}[flushleft]\footnotesize
\item Every verdict is generated from the estimates it describes; the sections give the designs and the caveats in full. Coefficients are in log points of the outcome per unit of the clustering share unless stated otherwise.
\end{tablenotes}
\end{threeparttable}
\end{table}


\subsection{Map of the document}

Section~\ref{sec:lit} says where each ingredient of the study comes from ---
every borrowed idea with its source. Section~\ref{sec:data} describes the feed,
the institutional details that govern the measures, and how the sample is
built, including when its filters bite. Section~\ref{sec:measures} is the
mathematical foundation: every measure's definition, the reference it
originates from, and the validation the implementation had to pass.
Section~\ref{sec:stylized} is the replication against the paper's published
magnitudes. Section~\ref{sec:liquidity} is the daily clustering--liquidity
relationship, Section~\ref{sec:book} the order-book anatomy and the within-day
results, and Section~\ref{sec:robust} the robustness work, including the
placebo, the sub-period splits, and an estimator cross-check that disagrees
informatively. Section~\ref{sec:strategy} prices a naive strategy on the
signal, in two variants. Section~\ref{sec:discussion} maps everything onto the
internship brief. The appendices carry the full variable definitions, the
ladder-inference algorithm with its hand-built test, and the reproduction
guide.
